\section{The trace format}\label{sec:format}

Section~\ref{sec:state} gave the trace as an abstract state \trace{}: an append-only log, a
lineage DAG of typed artifacts, a goal lens, a residual surface, and a derived freshness verdict.
This section makes that model concrete. The format is the load-bearing artifact of the whole
construction: the invariants of \S\ref{sec:state}, the certificate of \S\ref{sec:context}, and the
composition of \S\ref{sec:merge} are all statements \emph{about the typed objects defined here}. We
present the canonical types faithfully but distilled, and close with a worked interchange trace that
exhibits every discriminator the rest of the paper relies on.

\subsection{Canonical model and interchange projection}\label{sec:format-canonical}

The specification is layered. The \textbf{canonical model} is the semantic source of truth: strongly
typed, algebraic where a set of cases is closed, and the model against which invariants and semantics
are defined. The \textbf{interchange model} is a JSON projection of it, for persistence, transport,
and cross-language interoperability. The design rule is one-directional --- \emph{specify the
strongest semantic model first; derive the wire format from it, not the reverse} --- so the JSON is
never authoritative. This matters for verifiability: the invariants of \S\ref{sec:proved} are proved
over the canonical types, and the wire format inherits them only because it is a projection.

The canonical trace is a typed record whose fields are exactly the five parts of \S\ref{sec:state}
plus reproducibility and governance surfaces: an ordered \code{actions} list, a \code{meta\_actions}
overlay, an \code{artifacts} lineage DAG, \code{policies}/\code{policy\_evaluations},
\code{execution\_environments}, \code{goals}, and inter-trace \code{trace\_links}. Every algebraic
variant projects to JSON under an \textbf{explicit discriminator}: artifacts carry
\code{artifact\_type}, actions carry \code{category} and \code{type}, and verification-result variants
carry tags such as \code{proved}/\code{refuted}/\code{sat}/\code{unknown}. Decoding must reconstruct
the strict model and \emph{reject} an invalid discriminator/payload combination --- the wire boundary
is where ill-typed traces are caught, not admitted.

\subsection{Actions: the ground-truth log}\label{sec:format-actions}

Actions are the ordered, atomic steps --- one entry per tool call --- and they do not exchange
free-form names; they exchange \emph{artifact identities} via \code{inputs} and \code{outputs}. Every
action shares a common record (Figure~\ref{fig:fmt-action-common}).

\begin{figure}[H]
\centering
\[
\id{action\_common} \;::=\; \Big(\;
\begin{array}{@{}l@{\;:\;}l@{}}
\id{id}              & \id{Int},\\
\id{label}           & \id{String},\\
\id{rationale}       & \id{String},\\
\id{detail}          & \opt{\id{String}},\\
\id{inputs}          & \many{\id{ArtifactId}},\\
\id{outputs}         & \many{\id{ArtifactId}},\\
\id{evidence}        & \many{\id{Evidence}},\\
\id{observations}    & \many{\id{Observation}},\\
\id{execution}       & \opt{\id{ExecutionMetadata}},\\
\id{reproducibility} & \opt{\id{ActionReproducibility}},\\
\id{meta\_action\_id} & \opt{\id{String}}
\end{array}
\;\Big)
\]
\caption{The common record shared by every action: identity, labels, the \code{inputs}/\code{outputs} artifact-id edges, evidence and observations, and optional execution and reproducibility metadata.}
\label{fig:fmt-action-common}
\end{figure}

The strict action type partitions steps into four \textbf{categories}, each with its own closed set
of types: \emph{activity} (\code{ReadFile}, \code{EditFile}, \code{RunTests}, \code{GitCommit},
\code{FormulatePlan}, \ldots), \emph{gateway} (decisions, splits, loops --- carrying the options and
the reason each was chosen or rejected), \emph{reasoning} (\code{Formalize},
\code{DefineVerificationGoal}, \code{Verify}, \code{StateSpaceAnalysis}, \code{ConformanceCheck},
\code{CoSimulate}, \code{GenerateTests}), and \emph{governance} (\code{EvaluatePolicy},
\code{RequestApproval}, \code{RecordAudit}, \code{LinkTrace}, \ldots). The category/type split is the
semantic layer; the concrete engine is recorded separately in \code{execution\_metadata}
(\code{tool = "imandrax"}, \code{method\_ = "region\_decomposition"}, a \code{determinism} tag), which
keeps the format \emph{reasoner-agnostic} --- the same \code{StateSpaceAnalysis} action denotes the
same reasoning operation whether ImandraX or another engine realized it.

Over this atomic log sits an optional \textbf{meta-action} overlay: an ordered, non-overlapping
grouping of actions into units of \emph{intent} (a goal-to-steps hierarchy via \code{parent\_id}),
carrying a \code{title}, \code{intent}, \code{outcome}, and a \code{source} on a fidelity ladder
(\code{PlanDeclared} $>$ \code{TurnSegmented} $>$ \code{IntentInferred}). This is the \emph{curated
narrative} of a session: the atomic actions remain the evidence, and the overlay is a producer's
claim about how they group --- reviewable as a unit before drilling in.

\subsection{Artifacts and lineage}\label{sec:format-artifacts}

Artifacts are the evidence. Every artifact carries a common record (Figure~\ref{fig:fmt-artifact-common}) and then, for the structured
kinds, a typed payload aligned with its constructor.

\begin{figure}[H]
\centering
\small
\[
\id{artifact\_common} \;::=\; \Big(\;
\begin{array}{@{}l@{\;:\;}l@{}}
\id{artifact\_id}       & \id{String},\\
\id{artifact\_role}     & \opt{\id{ArtifactRole}},\\
\id{name}               & \opt{\id{String}},\\
\id{format}             & \opt{\id{String}},\\
\id{revision}           & \opt{\id{Int}},\\
\id{producer\_action\_id} & \opt{\id{Int}},\\
\id{derived\_from}      & \many{\id{ArtifactId}},\\
\id{supersedes}         & \opt{\id{String}},\\
\id{content\_ref}       & \opt{\id{String}},\\
\id{summary}            & \opt{\id{String}},\\
\id{component\_id}      & \opt{\id{String}},\\
\id{metadata}           & \opt{\id{ArtifactMetadata}}
\end{array}
\;\Big)
\]
\caption{The common record shared by every artifact: the record identity \code{artifact\_id}, semantic role, the \code{derived\_from} lineage edges, the \code{supersedes} revision link, and the rename-stable \code{component\_id}.}
\label{fig:fmt-artifact-common}
\end{figure}

\begin{sloppypar}
The artifact type is a \textbf{strict union} so that payload presence is aligned with kind and
impossible combinations are unrepresentable: source code, plans, diffs, and commits carry only the
common record; the reasoning kinds --- \code{FormalizationArtifact}, \code{FormalModelArtifact},
\code{VerificationGoalArtifact}, \code{VerificationResultArtifact},
\code{StateSpaceAnalysisResultArtifact}, \code{ConformanceResultArtifact},
\code{CoSimulationResultArtifact}, \code{GeneratedTestsArtifact} --- pair the common record with a
typed payload. Two members are newer additions the rest of the paper leans on: the residual surface
is folded in as a \code{Residual} artifact (\S\ref{sec:format-residual}), and \textbf{1.10 adds}
\code{CarriedForwardArtifact} --- the positive record of a merge skip (\S\ref{sec:format-carried}).
\end{sloppypar}

\code{derived\_from} \emph{is} the lineage DAG: each edge names an immediate input, so an artifact's
\textbf{dependency closure} is the transitive back-walk. This is what invariant I2 of
\S\ref{sec:state} constrains --- an edge points only at a strictly earlier artifact, forbidding
circular premises --- and what goal resolution and freshness walk. \textbf{Revisioning} is
append-only (I1): an artifact is immutable once produced, so re-analysis \emph{supersedes} (a new
\code{revision}, linked back by \code{supersedes}) rather than mutating in place. ``The latest
result'' is thus a derived query over the supersession chain, never an overwrite --- which is exactly
what lets an auditor see the whole history and what lets freshness be recomputed against any point.

\subsection{Formal results}\label{sec:format-results}

The reasoning payloads carry the object-logic verdicts. A \code{VerificationGoal} states what is to
be established (\code{kind} in verify/instance/theorem/lemma/axiom, a \code{target\_symbol}, a
list of \code{property\_item}s); a \code{VerificationResult} carries the outcome as a tagged variant
--- \code{ProvedResult} (with the pretty-printed proof), \code{RefutedResult} (with a
counterexample), \code{SatResult}, or \code{UnknownResult}. The state-space, conformance, and
co-simulation payloads are structurally analogous verdicts over their own object logics. Each result may
additionally carry a \code{reasoning\_fingerprint}, the anchor for freshness treated in its own right in
\S\ref{sec:freshness}.

\subsection{The residual surface: the negative space}\label{sec:format-residual}

A trace records what was established --- the positive space --- but review and agent-to-agent handoff
need the \emph{negative} space: what was \emph{not} established. The residual surface is the uniform,
typed, queryable record of it, and \textbf{a residual is an artifact}
(\code{artifact\_type = "Residual"}): it carries its fields in \code{payload} and \emph{anchors into
the lineage DAG} through \code{derived\_from}, so a gap hangs off exactly the artifact it qualifies.
Positive and negative space share one addressable model.

\begin{figure}[H]
\centering
\small
\[
\begin{array}{@{}l@{}}
\id{residual\_kind} \;::=\; \id{Assumption} \mid \id{Unverified} \mid \id{OutOfScope} \mid \id{Limitation} \mid \id{OpenQuestion} \mid{}\\
\hphantom{\id{residual\_kind} \;::=\;{}}\id{Defeater} \mid \id{NeedsRereasoning} \mid \id{CoverageRegression}
\end{array}
\]
\begin{center}
\begin{tabular}{@{}ll@{}}
$\id{Assumption}$         & \text{\footnotesize\itshape a premise relied on but not established}\\
$\id{Unverified}$         & \text{\footnotesize\itshape something done but not checked or proved}\\
$\id{OutOfScope}$         & \text{\footnotesize\itshape deliberately not addressed}\\
$\id{Limitation}$         & \text{\footnotesize\itshape a constraint under which the results hold}\\
$\id{OpenQuestion}$       & \text{\footnotesize\itshape a decision deferred to review}\\
$\id{Defeater}$           & \text{\footnotesize\itshape counter-evidence against a claim (not a gap); carries a $\id{defeater\_kind}$}\\
$\id{NeedsRereasoning}$   & \text{\footnotesize\itshape a result a merge disturbed; re-establish against the merged code}\\
$\id{CoverageRegression}$ & \text{\footnotesize\itshape a goal whose scope gained an unproven member after a merge}
\end{tabular}
\end{center}
\[
\id{defeater\_kind} \;::=\; \id{Rebuts} \mid \id{Undermines} \mid \id{Undercuts}
\]
\begin{center}
\begin{tabular}{@{}ll@{}}
$\id{Rebuts}$     & \text{\footnotesize\itshape the claim may be false --- e.g.\ a counterexample to a proved property}\\
$\id{Undermines}$ & \text{\footnotesize\itshape the evidence is invalid --- e.g.\ the model does not match the code}\\
$\id{Undercuts}$  & \text{\footnotesize\itshape the inference is deficient --- e.g.\ ``passing tests do not establish this''}
\end{tabular}
\end{center}
\caption{The residual taxonomy (Pollock / SEI eliminative-argumentation): \code{residual\_kind} enumerates the gap and counter-evidence kinds (including the merge-specific \code{NeedsRereasoning} and \code{CoverageRegression}), and \code{defeater\_kind} classifies a defeater as a rebut, undermine, or undercut.}
\label{fig:fmt-residual-kind}
\end{figure}

The first five kinds record \emph{missing positive space} --- something not established. A
\code{Defeater} is different in kind: it records \emph{negative evidence}, a concrete reason to
believe a stated result is \emph{wrong}, and it anchors (via \code{derived\_from}/\code{target}) to
the claim it contests while citing the counter-evidence in \code{related\_artifact\_ids}. An
\emph{open} defeater makes its target \textbf{contested} --- a \code{Property} acceptance item over a
contested proof resolves \code{AcceptBlocked}, not \code{AcceptDone} (this is the non-monotonic
retraction of \S\ref{sec:context}, made first-class and reviewer-raisable). \textbf{1.10 adds} two
merge-specific kinds: \code{NeedsRereasoning} (an established result whose evidence a merge disturbed)
and \code{CoverageRegression} (a goal whose scope gained an unproven member) --- the negative half of
the composition account of \S\ref{sec:merge}. Each residual also carries a \code{severity}
(impact-if-wrong, not probability), a \code{source} on a trust ladder (\code{AgentDeclared},
\code{PolicyDerived}, \code{ToolInferred}, \code{ReviewerAdded}), a \code{status}
(\code{Open}/\code{Acknowledged}/\code{Addressed}/\code{Waived}), and an optional
\code{suggested\_check} that turns a warning into a work-item. Because it is an ordinary artifact,
every severity/status/source field and every policy applies to a defeater unchanged.

\subsection{Carried-forward: the positive dual}\label{sec:format-carried}

Composition (\S\ref{sec:merge}) needs a symmetric record on the positive side. \textbf{1.10 adds} the
\code{CarriedForwardArtifact}: where a residual records a gap, a \code{CarriedForward} records a claim
that was \emph{deliberately not re-checked because it is provably safe}. It carries a falsifiable
witness so an auditor can re-verify the skip (Figure~\ref{fig:fmt-carried-forward}).

\begin{figure}[H]
\centering
\small
\[
\id{carried\_forward\_basis} \;::=\; \id{ClosureDisjoint} \mid \id{UninterpretedOpaque}
\]
\begin{center}
\begin{tabular}{@{}ll@{}}
$\id{ClosureDisjoint}$     & \text{\footnotesize\itshape the result's dependency closure did not intersect the merge's change set}\\
$\id{UninterpretedOpaque}$ & \text{\footnotesize\itshape every touched dependency was assumed uninterpreted}
\end{tabular}
\end{center}
\[
\id{carried\_forward\_payload} \;::=\; \Big(\;
\begin{array}{@{}l@{\;:\;}l@{}}
\id{result\_id}       & \id{String},\\
\id{basis}            & \id{carried\_forward\_basis},\\
\id{closure}          & \many{\id{ComponentId}},\\
\id{via\_assumptions} & \many{\id{String}}
\end{array}
\;\Big)
\]
\caption{The \code{CarriedForward} record documenting a provably safe merge skip: the \code{basis} (disjoint closure or uninterpreted dependency) and the \code{closure} of component ids that serves as the auditor's falsifiable witness.}
\label{fig:fmt-carried-forward}
\end{figure}

The merged trace's account of every prior result is thus exactly one of three under the
\textbf{totality} invariant of \S\ref{sec:merge}: \code{CarriedForward}, \code{NeedsRereasoning}, or
freshly re-established. The two-parent \code{merge\_event} that occasions this is recorded on the
merged trace as provenance kept \emph{off} the artifact DAG --- so I2 acyclicity (a \code{derived\_from}
edge points only at an earlier artifact) survives composition. The change set is matched by
\code{component\_id}, so a rename is one component rather than a delete-plus-add; detail is deferred to
\S\ref{sec:merge}.

\subsection{Reproducibility and execution environments}\label{sec:format-repro}

Reproducibility is recorded at two granularities. At the \textbf{trace} level, a
\code{trace\_reproducibility} record states a \code{status}
(\code{NotReproducible}/\code{Partially}/\code{Reproducible}), entrypoints, and the required artifacts
and environments. At the \textbf{action} level, each consequential action may carry a
\code{reproduction\_kind} (deterministic replay, tool re-execution, procedural or manual replay), a
concrete \code{procedure} (command, working directory, arguments, or steps), and an
\code{expected\_output}. Both refer to named \code{execution\_environment}s --- a toolchain,
container, service, or external system with its versioned components --- so a reader can reconstruct
\emph{where} a result was produced, not merely that it was. The discipline is that a consequential
reasoning outcome is either directly reproducible or linked to a reproducible downstream validation.

\subsection{A worked trace}\label{sec:format-example}

Figure~\ref{fig:fmt-worked-trace} shows the interchange projection of a minimal but complete slice: one reasoning
action, the \code{VerificationGoal} it consumed and the \code{VerificationResult} it produced (with a
\code{reasoning\_fingerprint}), and one \code{Residual} contesting that result. It exhibits the
\code{category}/\code{artifact\_type} discriminators and the \code{derived\_from} lineage that
carries the whole construction.

\begin{figure}[H]
\begin{lstlisting}[language=,basicstyle=\footnotesize\ttfamily]
{
  "actions": [
    { "id": 4, "category": "reasoning", "type": "Verify",
      "label": "Verify amount_inv on the payment model",
      "rationale": "capture must never exceed the authorized amount",
      "inputs": ["a2", "g1"], "outputs": ["vr1"],
      "execution": { "tool": "imandrax", "version": "2026.04",
                     "method": "model_check", "determinism": "deterministic" } }
  ],
  "artifacts": [
    { "artifact_id": "g1", "artifact_type": "VerificationGoal",
      "producer_action_id": 3, "derived_from": ["a2"],
      "component_id": "cmp:amount_inv",
      "payload": { "goal_id": 1, "kind": "VerifyGoal",
                   "description": "amount_captured <= amount",
                   "target_artifact_id": "a2", "target_symbol": "amount_inv",
                   "properties": [ { "name": "amount_inv", "status": "Pending",
                                     "src": "amount_captured <= amount" } ] } },

    { "artifact_id": "vr1", "artifact_type": "VerificationResult",
      "artifact_role": "proof", "producer_action_id": 4,
      "derived_from": ["g1", "a2"], "component_id": "cmp:amount_inv",
      "summary": "amount_inv proved",
      "payload": { "goal_id": 1, "goal_artifact_id": "g1", "status": "proved",
                   "engine": "imandrax",
                   "result": { "proved": { "proof_pp": "...", "properties": [] } },
                   "fingerprint": {
                     "task_checksum": "sha256:7f3c...e1",   /* target + closure */
                     "task_shape": "sha256:0a9b...",
                     "target_symbol": "amount_inv",
                     "engine": "imandrax", "engine_version": "2026.04",
                     "model_artifact_id": "a2", "model_revision": 1 } } },

    { "artifact_id": "r1", "artifact_type": "Residual",
      "derived_from": ["vr1", "a13"],           /* anchors to the claim it contests */
      "summary": "amount_inv refuted under interleaved capture",
      "payload": { "kind": "defeater", "defeater_kind": "rebuts",
                   "statement": "capture can exceed amount under concurrent capture/refund",
                   "severity": "critical",
                   "target": { "target_type": "artifact", "target_id": "vr1" },
                   "related_artifact_ids": ["a13"],   /* the counterexample */
                   "source": "reviewer_added", "status": "open" } }
  ]
}
\end{lstlisting}
\caption{A worked interchange trace: a \code{Verify} reasoning action, the \code{VerificationGoal} \code{g1} and proved \code{VerificationResult} \code{vr1} (with its \code{reasoning\_fingerprint}), and an open \code{Residual} defeater \code{r1} that rebuts \code{vr1} --- exhibiting the discriminators and \code{derived\_from} lineage the rest of the paper relies on.}
\label{fig:fmt-worked-trace}
\end{figure}

Reading it as a context: \code{vr1} is a \emph{proved} result, and its fingerprint lets a consumer
recompute whether it is still \code{Fresh} against the current model. But \code{r1} is an
\emph{open} \code{Defeater} that \code{Rebuts} it (anchored by \code{derived\_from} and \code{target}),
so the claim is \emph{contested}: a \code{Property} obligation over \code{vr1} resolves
\code{AcceptBlocked}, not \code{AcceptDone}. The certificate travels with the claim --- which oracle
(\code{imandrax}), under what task (the fingerprinted closure), whether fresh, whether contested ---
which is precisely what the freshness of \S\ref{sec:proved} and the composition of \S\ref{sec:merge}
compute over.
